%=============================================================================
% WAVE 3, ITERATION 7: PATENT 6 + PATENT 8 + PATENT 9 COMBINED UPDATE
% Decihertz Detector Comparison, Resonant Cavity Transmitters, SQL Positioning
% Agent 6/8/9 (W3i7)
% Date: 20 March 2026
%
% Builds on:
%   W3i6_P6_P7_P9_update.tex (P6 all confirmed; dpsiem noise 22 orders below)
%   W3i6_P8_P11_update.tex   (P8 submarine killer app; 10^4x mass gap; DSN 370x)
%   W3i5_Cosmo_MatSci_update.tex (two-component decomposition)
%   W3i4_P6_update.tex       (P6 sensitivity: h_min = 9.3e-16/sqrt(Hz) at N=10)
%   W3i4_P9_update.tex       (P9 positioning: 8.7 nm Heisenberg-limited)
%
% W3i7 MANDATE:
%   P6: Complete comparison table vs LISA, DECIGO, BBO, ET, CE for 0.001-10 Hz.
%       Where is DFD P6 uniquely competitive?
%   P8: Resonant cavity-based psi transmitters with Q_psi amplification?
%       Minimum transmitter mass for 1 kbit/s at 1 AU.
%   P9: Realistic (SQL) performance vs Heisenberg limit.
%       What applications survive at SQL (8.7 um) vs Heisenberg (8.7 nm)?
%
% KEY W3i6 INPUTS:
%   1. P6: h_min(0.1 Hz, N=10, K=3) = 9.3e-16 /sqrt(Hz), RPN-limited
%   2. P6: Coherent array h_min ~ 10^-20 to 10^-21 at N=10^5-10^6
%   3. P8: Submarine killer app; M_s >= 10^10 kg for 1 kbit/s at 500 m
%   4. P8: 10^4x mass gap between achievable (~10^6 kg) and required (~10^10 kg)
%   5. P9: 8.7 nm positioning (Heisenberg-limited); 0.091 mm single-sensor
%   6. P9: Geological signal/EM noise separation > 10^25
%=============================================================================

\documentclass[11pt]{article}
\usepackage[margin=2cm]{geometry}
\usepackage{amsmath,amssymb,amsthm,mathtools}
\usepackage{physics}
\usepackage{booktabs}
\usepackage{longtable}
\usepackage{hyperref}
\usepackage{xcolor}
\usepackage{array}
\usepackage{siunitx}
\usepackage{tcolorbox}
\usepackage{enumitem}
\usepackage{float}

\newtheorem{theorem}{Theorem}[section]
\newtheorem{proposition}[theorem]{Proposition}
\newtheorem{corollary}[theorem]{Corollary}
\newtheorem{lemma}[theorem]{Lemma}
\theoremstyle{definition}
\newtheorem{definition}[theorem]{Definition}
\theoremstyle{remark}
\newtheorem{remark}[theorem]{Remark}
\newtheorem{finding}[theorem]{Finding}
\newtheorem{correction}[theorem]{Correction}
\newtheorem{result}[theorem]{Result}

\newcommand{\ee}[1]{\times10^{#1}}
\newcommand{\lambdaone}{|\lambda-1|}
\newcommand{\Meff}{M_{\mathrm{eff}}}
\newcommand{\Mpsi}{M_\psi}
\newcommand{\Qpsi}{Q_\psi}
\newcommand{\Opsi}{\Omega_\psi}
\newcommand{\sqrtHz}{\,\mathrm{Hz}^{-1/2}}
\newcommand{\SNR}{\mathrm{SNR}}
\newcommand{\Msun}{M_\odot}
\newcommand{\kB}{k_{\mathrm{B}}}
\newcommand{\Pmax}{P_{\mathrm{max}}}
\newcommand{\psigrav}{\psi_{\mathrm{grav}}}
\newcommand{\dpsiem}{\delta\psi_{\mathrm{EM}}}
\newcommand{\GN}{G_N}

\newcommand{\confirmed}[1]{\textcolor{blue}{\textbf{CONFIRMED:} #1}}
\newcommand{\corrected}[1]{\textcolor{red}{\textbf{CORRECTED:} #1}}
\newcommand{\caution}[1]{\textcolor{orange}{\textbf{CAUTION:} #1}}
\newcommand{\honest}[1]{\textcolor{violet}{\textbf{HONEST ASSESSMENT:} #1}}
\newcommand{\killed}[1]{\textcolor{red!70!black}{\textbf{KILLED:} #1}}
\newcommand{\verdict}[1]{\textcolor{green!50!black}{\textbf{VERDICT:} #1}}

\tcbuselibrary{skins,breakable}

\title{\textbf{Wave 3 Iteration 7: Patent 6 + Patent 8 + Patent 9 Update}\\[6pt]
Decihertz GW Detection: Complete Comparison Table (P6),\\
Resonant Cavity Psi-Transmitters and Minimum Mass (P8),\\
SQL vs Heisenberg Positioning: Application Survival (P9)}
\author{DFD Research Programme --- Agent 6/8/9, W3i7}
\date{20 March 2026}

\begin{document}
\maketitle
\tableofcontents
\newpage

%=============================================================================
\section*{Executive Summary}
%=============================================================================

This document addresses three mandates from the W3i7 task specification:

\begin{enumerate}
\item \textbf{P6 (Decihertz GW Detection --- Complete Comparison):}
  A comprehensive frequency-band-by-frequency-band comparison of the DFD P6
  atom-interferometer array against LISA, DECIGO, BBO, Einstein Telescope (ET),
  and Cosmic Explorer (CE). The DFD P6 system is \textbf{uniquely competitive
  in the 0.01--1~Hz ``decihertz gap''} at array sizes $N \geq 10^4$, where no
  funded detector currently operates. Below 0.01~Hz, LISA dominates. Above 1~Hz,
  ET and CE dominate. In the decihertz band, DECIGO and BBO are the only
  competitors, but both are unfunded concept studies at TRL~2--3, while DFD P6
  builds on demonstrated atom-interferometry technology (MAGIS, AION, ZAIGA)
  at TRL~3--4.

\item \textbf{P8 (Resonant Cavity Psi-Transmitters):}
  We investigate whether $\Qpsi$ amplification in an SRF cavity can reduce the
  required transmitter mass. The answer is \textbf{no}: the $\Qpsi$ factor
  amplifies the \emph{stored energy} in the cavity mode, but the far-field
  psi-radiation scales as the \emph{power radiated}, not the stored energy.
  The radiated power is $P_{\mathrm{rad}} = \Opsi E_{\mathrm{stored}}/\Qpsi$,
  which is \emph{independent} of $\Qpsi$ at fixed input power. The minimum
  transmitter mass for 1~kbit/s at 1~AU is $M_s \gtrsim 2.7\ee{17}$~kg
  ($\sim 4.5\ee{-8}\,\Msun$, roughly Ceres-mass). The mass gap from any
  foreseeable technology is $\sim 10^{11}\times$.

\item \textbf{P9 (SQL vs Heisenberg Positioning):}
  The 8.7~nm positioning claim assumes Heisenberg-limited (HL) sensors.
  At the standard quantum limit (SQL), the positioning precision degrades
  by $\sqrt{N_{\mathrm{photon}}}$ for optical readout or $\sqrt{N_{\mathrm{atom}}}$
  for atom interferometry: the SQL array precision is
  $\sigma_r^{\mathrm{SQL}} \approx 8.7~\mu\mathrm{m}$, a factor of $10^3$
  worse than HL. Applications that survive at SQL: submarine navigation
  ($<1$~m required), geological survey (sub-m adequate), seismic monitoring,
  autonomous vehicle navigation, and infrastructure monitoring. Applications
  that \emph{require} HL: chip-scale metrology, interferometric geodesy at
  nm-level, and gravitational-wave detector calibration.
\end{enumerate}

\begin{tcolorbox}[colback=blue!5!white, colframe=blue!60!black,
  title=\textbf{CRITICAL W3i7 RESULTS --- READ FIRST}]
\begin{enumerate}[nosep]
\item \textbf{P6}: DFD P6 is uniquely competitive in the 0.01--1~Hz decihertz
  gap at $N \geq 10^4$ sensors. No funded detector covers this band.
\item \textbf{P8}: $\Qpsi$ amplification does NOT reduce transmitter mass.
  Minimum mass for 1~kbit/s at 1~AU: $2.7\ee{17}$~kg (Ceres-mass).
  Psi-comms remains technologically infeasible by $\sim 10^{11}\times$.
\item \textbf{P9}: SQL positioning is $8.7~\mu$m (not 8.7~nm).
  Most practical applications survive at SQL. The 8.7~nm HL claim is
  aspirational but not required for the dominant use cases.
\end{enumerate}
\end{tcolorbox}

%=============================================================================
\part{Patent 6: Decihertz GW Detection --- Complete Detector Comparison}
\label{part:P6}
%=============================================================================

%=============================================================================
\section{The Gravitational Wave Detection Landscape: 0.001--10~Hz}
\label{sec:P6-landscape}
%=============================================================================

\subsection{Existing and Proposed Detectors}

The gravitational wave frequency spectrum from $10^{-4}$~Hz to $10^4$~Hz is
covered (or proposed to be covered) by five classes of detector:

\begin{enumerate}[nosep]
\item \textbf{LISA} (Laser Interferometer Space Antenna):
  Space-based, $2.5\ee{9}$~m arms, sensitive from $\sim 10^{-4}$~Hz to
  $\sim 0.1$~Hz. Funded (ESA/NASA), launch $\sim$2035. Design sensitivity
  $h_{\min} \sim 10^{-20}\sqrtHz$ at $f \sim 3$~mHz (optimal).

\item \textbf{DECIGO} (DECi-hertz Interferometer Gravitational wave
  Observatory): Space-based, 1000~km arms, Fabry--Perot cavities.
  Sensitive from 0.1~Hz to 10~Hz. \emph{Unfunded concept study} (JAXA).
  Projected sensitivity $h_{\min} \sim 10^{-23}\sqrtHz$ at 0.1--1~Hz.

\item \textbf{BBO} (Big Bang Observer): Space-based, similar to DECIGO but
  with US origin. Sensitive from 0.01~Hz to 10~Hz. \emph{Unfunded concept
  study} (NASA). Projected sensitivity $h_{\min} \sim 4\ee{-24}\sqrtHz$
  at 0.3~Hz.

\item \textbf{ET} (Einstein Telescope): Underground, triangular,
  10~km arms, cryogenic. Sensitive from $\sim 2$~Hz to $10^4$~Hz.
  Funded (EU), construction $\sim$2030s. Design sensitivity
  $h_{\min} \sim 10^{-24}\sqrtHz$ at 3--10~Hz (low-frequency corner
  limited by Newtonian noise).

\item \textbf{CE} (Cosmic Explorer): Surface, L-shaped, 40~km arms.
  Sensitive from $\sim 5$~Hz to $10^4$~Hz. \emph{Proposed} (US NSF),
  partially funded for design study. Design sensitivity
  $h_{\min} \sim 3\ee{-25}\sqrtHz$ at 10--100~Hz.
\end{enumerate}

\subsection{The Decihertz Gap}

The frequency range 0.01--1~Hz is often called the ``decihertz gap'' because:
\begin{itemize}[nosep]
\item LISA sensitivity degrades above $\sim 0.03$~Hz (arm length too short).
\item ET/CE are limited below $\sim 2$~Hz by seismic and Newtonian noise.
\item DECIGO and BBO would fill this gap but neither is funded.
\end{itemize}

Ground-based atom interferometers (MAGIS-100, AION-10, ZAIGA) target
$0.01$--$10$~Hz but achieve $h_{\min} \sim 10^{-15}$--$10^{-17}\sqrtHz$
in single-baseline configurations --- insufficient for astrophysical sources.

\subsection{DFD P6 Array Sensitivity}

From W3i4 and W3i6 (confirmed), the DFD P6 atom interferometer sensitivity
with $N$ sensors and $K = 3$ resonant enhancement is:
\begin{equation}
h_{\min}(f, N) = \frac{h_1(f)}{\sqrt{N}},
\label{eq:P6-scaling}
\end{equation}
where the single-sensor RPN-limited strain noise at optimal frequency
$f \approx 0.1$~Hz is:
\begin{equation}
h_1(0.1\,\mathrm{Hz}) = 9.3\ee{-16} \times \sqrt{10}\sqrtHz
= 2.94\ee{-15}\sqrtHz.
\label{eq:h1-single}
\end{equation}
(The W3i4 value $9.3\ee{-16}\sqrtHz$ is for $N = 10$; the single-sensor
value is $\sqrt{10}$ times worse.)

The frequency dependence of the atom interferometer transfer function is:
\begin{equation}
\mathcal{T}(f, T) = \sin^2(\pi f T) \cdot \frac{1}{(\pi f T)^2}
\cdot (4\pi f T)^2 = 16\sin^2(\pi f T),
\label{eq:transfer}
\end{equation}
with interrogation time $T = 1$~s (baseline). The sensitivity peaks at
$f = 1/(2T) = 0.5$~Hz and has nulls at $f = n/T$. For broadband analysis
we use the envelope $\mathcal{T} \sim \min(1, (f/f_0)^2)$ below the peak
and $\mathcal{T} \sim 1$ near the peak.

The full strain sensitivity for an $N$-sensor array:
\begin{equation}
h_{\min}(f, N) \approx \frac{2.94\ee{-15}}{\sqrt{N}}
\cdot \frac{1}{\mathcal{T}(f, 1\,\mathrm{s})}\sqrtHz.
\label{eq:P6-full}
\end{equation}

%=============================================================================
\section{Complete Comparison Table: 0.001--10~Hz}
\label{sec:P6-comparison}
%=============================================================================

\begin{table}[H]
\centering
\caption{Gravitational wave detector strain sensitivity $h_{\min}$
($\mathrm{Hz}^{-1/2}$) across the 0.001--10~Hz band. DFD~P6 shown for
$N = 10^2, 10^4, 10^6$ sensors. Bold indicates best detector in each band.
``---'' means detector has no useful sensitivity in that band.
$\dagger$~=~unfunded concept.}
\label{tab:P6-comparison}
\small
\begin{tabular}{@{}lccccccccc@{}}
\toprule
$f$ (Hz) & LISA & DECIGO$^\dagger$ & BBO$^\dagger$ & ET & CE
  & DFD $10^2$ & DFD $10^4$ & DFD $10^6$ & Winner \\
\midrule
0.001 & $\mathbf{4\ee{-20}}$ & --- & --- & --- & ---
  & $3\ee{-8}$ & $3\ee{-9}$ & $3\ee{-10}$
  & \textbf{LISA} \\[3pt]
0.003 & $\mathbf{1\ee{-20}}$ & --- & $2\ee{-22}$ & --- & ---
  & $3\ee{-10}$ & $3\ee{-11}$ & $3\ee{-12}$
  & \textbf{LISA} \\[3pt]
0.01 & $\mathbf{4\ee{-20}}$ & $5\ee{-22}$ & $1\ee{-22}$ & --- & ---
  & $3\ee{-12}$ & $3\ee{-13}$ & $3\ee{-14}$
  & \textbf{LISA} \\[3pt]
0.03 & $3\ee{-18}$ & $\mathbf{2\ee{-23}}$ & $\mathbf{8\ee{-24}}$ & --- & ---
  & $3\ee{-13}$ & $3\ee{-14}$ & $3\ee{-15}$
  & \textbf{BBO$^\dagger$} \\[3pt]
0.1 & $2\ee{-17}$ & $\mathbf{1\ee{-23}}$ & $\mathbf{4\ee{-24}}$ & --- & ---
  & $2.9\ee{-14}$ & $2.9\ee{-15}$ & $\mathbf{2.9\ee{-16}}$
  & BBO$^\dagger$ \\[3pt]
0.3 & --- & $\mathbf{1\ee{-23}}$ & $\mathbf{4\ee{-24}}$ & --- & ---
  & $3\ee{-14}$ & $3\ee{-15}$ & $3\ee{-16}$
  & BBO$^\dagger$ \\[3pt]
1.0 & --- & $\mathbf{2\ee{-23}}$ & $\mathbf{5\ee{-24}}$ & $5\ee{-22}$ & ---
  & $3\ee{-14}$ & $3\ee{-15}$ & $3\ee{-16}$
  & BBO$^\dagger$ \\[3pt]
3.0 & --- & $1\ee{-22}$ & $3\ee{-23}$ & $\mathbf{1\ee{-24}}$ & $5\ee{-24}$
  & $8\ee{-14}$ & $8\ee{-15}$ & $8\ee{-16}$
  & \textbf{ET} \\[3pt]
10 & --- & --- & --- & $\mathbf{3\ee{-25}}$ & $\mathbf{3\ee{-25}}$
  & --- & --- & ---
  & \textbf{ET/CE} \\
\bottomrule
\end{tabular}
\end{table}

\subsection{Key Observations from the Comparison Table}

\begin{finding}[DFD P6 Absolute Sensitivity vs Dedicated GW Detectors]
\label{find:P6-absolute}
At all frequencies and array sizes considered, the DFD P6 atom-interferometer
array is \textbf{less sensitive} than the dedicated GW detector concepts
(LISA, DECIGO, BBO, ET, CE) by margins ranging from $10^{6}\times$ to
$10^{10}\times$ in strain.

\textbf{Root cause}: The DFD P6 sensitivity is limited by the radiation
pressure noise (RPN) of the atom interferometer at
$h_1 \sim 3\ee{-15}\sqrtHz$. Even with $N = 10^6$ sensors ($\sqrt{N}$
improvement of $10^3$), the best achievable is $h_{\min} \sim 3\ee{-16}\sqrtHz$.
LISA achieves $\sim 10^{-20}\sqrtHz$ with a single 2.5-Gm baseline. The
$10^4\times$ deficit reflects the fundamental difference between
km/Gm-scale laser interferometry and $\sim$m-scale atom interferometry.
\end{finding}

\begin{finding}[DFD P6 Cannot Compete on Raw Sensitivity]
\label{find:P6-nocompete}
\killed{DFD P6 does NOT compete with LISA, DECIGO, BBO, ET, or CE on
raw strain sensitivity at any frequency.}

The W3i6 statement that P6 is ``LISA-competitive at $N \geq 3.2\ee{3}$''
requires careful qualification. This referred to reaching LISA's sensitivity
\emph{at LISA's worst frequency} ($\sim 0.1$~Hz, where LISA performance
degrades to $h \sim 10^{-17}$). At LISA's \emph{optimal} frequency
($\sim 3$~mHz), LISA outperforms any conceivable DFD P6 array by $> 10^6\times$.
\end{finding}

%=============================================================================
\section{Where Is DFD P6 Uniquely Competitive?}
\label{sec:P6-unique}
%=============================================================================

\subsection{The Funded-Detector Argument}

While DFD P6 loses on absolute sensitivity to DECIGO and BBO, those
detectors are \emph{unfunded concept studies}:

\begin{table}[H]
\centering
\caption{Technology readiness and funding status of decihertz detectors.}
\label{tab:P6-TRL}
\begin{tabular}{@{}lccl@{}}
\toprule
Detector & TRL & Funding Status & Timeline \\
\midrule
LISA & 5--6 & Funded (ESA L3) & Launch $\sim$2035 \\
DECIGO & 2--3 & Unfunded; pathfinder (B-DECIGO) proposed & $>$2040 \\
BBO & 2 & Unfunded concept & $>$2045 \\
ET & 3--4 & Funded (EU, site selected) & First light $\sim$2035 \\
CE & 2--3 & Design study funded (NSF) & $>$2040 \\
\midrule
MAGIS-100 & 3--4 & Funded (DOE) & Operating $\sim$2027 \\
AION-10 & 3 & Funded (UKRI) & Operating $\sim$2026 \\
DFD P6 ($N=10$) & 2--3 & Concept & --- \\
DFD P6 ($N=10^4$) & 1--2 & Concept & --- \\
\bottomrule
\end{tabular}
\end{table}

\begin{finding}[DFD P6 Unique Niche: Decihertz Monitoring with Existing Technology]
\label{find:P6-niche}
\verdict{DFD P6 occupies a unique niche NOT because it beats other detectors
on sensitivity, but because:}
\begin{enumerate}[nosep]
\item \textbf{No funded detector covers 0.01--1~Hz.} LISA cuts off above
  $\sim 0.03$~Hz; ET/CE cut off below $\sim 2$~Hz. The decihertz gap will
  persist until DECIGO/BBO are funded (timeline $>$2040).
\item \textbf{Atom interferometry is ground-based.} Unlike DECIGO/BBO, DFD P6
  does not require a space mission. It leverages MAGIS/AION technology.
\item \textbf{Scalability through array size.} Adding sensors improves
  sensitivity as $1/\sqrt{N}$, allowing incremental deployment.
\item \textbf{Multi-messenger astronomy.} Even at $h \sim 10^{-15}$, a
  decihertz detector can observe the \emph{inspiral} phase of stellar-mass
  binary black holes days before merger (when the signal sweeps through
  $\sim 0.1$~Hz at $h \sim 10^{-18}$ for sources at $\sim 100$~Mpc).
\end{enumerate}
\end{finding}

\subsection{What Sources Can DFD P6 Detect?}

\begin{table}[H]
\centering
\caption{Astrophysical sources in the decihertz band and DFD P6 detectability.
SNR $>$ 5 required for detection.}
\label{tab:P6-sources}
\begin{tabular}{@{}p{3.5cm}ccccl@{}}
\toprule
Source & $f$ (Hz) & $h$ at $d$ & $d$ & DFD $N$ for SNR $= 5$
  & Notes \\
\midrule
IMBH binary ($10^3$--$10^4\,\Msun$) merger
  & 0.1--1 & $10^{-18}$ & 1~Gpc
  & $2.4\ee{6}$ & W3i4 confirmed \\[3pt]
Stellar BBH inspiral (30+30~$\Msun$)
  & 0.03--0.3 & $5\ee{-20}$ & 100~Mpc
  & $8.6\ee{9}$ & Impractical \\[3pt]
Stellar BBH inspiral (30+30~$\Msun$)
  & 0.1 & $10^{-17}$ & 10~Mpc
  & $8.6\ee{5}$ & Rare sources \\[3pt]
SGWB ($\Omega_{\rm GW} = 10^{-9}$)
  & 0.01--1 & $\sim 10^{-22}$ & ---
  & $\sim 10^{14}$ & Impractical \\[3pt]
WD binary (verification)
  & 0.001--0.01 & $10^{-21}$ & kpc
  & $>10^{12}$ & LISA domain \\
\bottomrule
\end{tabular}
\end{table}

\begin{finding}[DFD P6 Realistic Detection Targets]
\label{find:P6-targets}
The only astrophysical source detectable by a DFD P6 array of feasible
size ($N \leq 10^6$) is the \textbf{IMBH binary merger} at
$10^3$--$10^4\,\Msun$ in the frequency range 0.1--1~Hz. This requires
$N \approx 2.4\ee{6}$ sensors for SNR~$= 5$ at 1~Gpc.

\honest{Stellar-mass BBH inspirals and the stochastic GW background
require $N > 10^9$, which is impractical. DFD P6 is a
\emph{single-source-class} decihertz detector, not a general-purpose
GW observatory.}
\end{finding}

%=============================================================================
\section{DFD P6 vs Competitors: Definitive Band-by-Band Verdict}
\label{sec:P6-verdict}
%=============================================================================

\begin{table}[H]
\centering
\caption{Band-by-band verdict: which detector wins in each frequency range?}
\label{tab:P6-band-verdict}
\begin{tabular}{@{}p{2.5cm}p{3cm}p{3cm}p{5cm}@{}}
\toprule
Band (Hz) & Winner (Funded) & Winner (All) & DFD P6 Status \\
\midrule
$10^{-4}$--$10^{-2}$
  & LISA & LISA
  & No competition; $>10^{6}\times$ worse \\[4pt]
$10^{-2}$--$10^{-1}$
  & \textbf{None funded} & BBO/DECIGO
  & Unique ground-based option; $10^{7}\times$ worse than BBO
    but BBO is unfunded \\[4pt]
$10^{-1}$--$10^{0}$
  & \textbf{None funded} & BBO/DECIGO
  & \textbf{Strongest relative position}; only ground-based
    decihertz option with demonstrated technology base \\[4pt]
$10^{0}$--$10^{1}$
  & ET & ET
  & No competition; ET $>10^{8}\times$ better \\[4pt]
$> 10$
  & ET/CE & CE
  & Not applicable (AI transfer function nulls) \\
\bottomrule
\end{tabular}
\end{table}

\begin{tcolorbox}[colback=green!5!white, colframe=green!40!black,
  title=\textbf{P6 Decihertz Verdict: Strongest Claim is Technology Accessibility}]
DFD P6 does \textbf{not} win on raw sensitivity in any band. Its unique
value proposition is:
\begin{enumerate}[nosep]
\item \textbf{Ground-based decihertz access} when no funded space mission
  covers 0.01--1~Hz.
\item \textbf{Incremental scalability} --- deploy 10 sensors, then 100,
  then 1000, characterising noise at each stage.
\item \textbf{IMBH detection} at $N \sim 10^6$ in a band inaccessible to
  LISA or ET/CE.
\item \textbf{Technology synergy} with MAGIS/AION/ZAIGA atom interferometry
  programmes already funded and under construction.
\end{enumerate}
The patent claim of ``decihertz GW detection'' is \textbf{valid but
context-dependent}: it fills a gap left by the absence of funded
space-based decihertz missions, not by outperforming them.
\end{tcolorbox}

%=============================================================================
\part{Patent 8: Resonant Cavity Psi-Transmitters and Minimum Mass}
\label{part:P8}
%=============================================================================

%=============================================================================
\section{Can $\Qpsi$ Amplification Reduce Transmitter Mass?}
\label{sec:P8-Qpsi}
%=============================================================================

\subsection{The Idea: Coherent Psi-Field Buildup in an SRF Cavity}

The proposal is to use a high-$Q$ SRF cavity (as in P7/P11) to build up a
coherent psi-field mode, then radiate it as a communication signal. The
stored psi-energy in a cavity of quality factor $\Qpsi$ and input power
$P_{\mathrm{in}}$:
\begin{equation}
E_{\mathrm{stored}} = \frac{\Qpsi}{\Opsi} P_{\mathrm{in}}
= \tau_{1/e} \cdot P_{\mathrm{in}}.
\label{eq:E-stored}
\end{equation}
For $\Qpsi = 10^9$, $\Opsi = 10^9$~rad/s: $E_{\mathrm{stored}} = P_{\mathrm{in}}$
(in joules when $P_{\mathrm{in}}$ is in watts). The cavity stores 1~second's
worth of input power.

The hope is that the $\Qpsi = 10^9$ amplification of the intracavity
psi-field amplitude ($\psi_{\mathrm{cav}} \sim \sqrt{\Qpsi}\,\psi_{\mathrm{in}}$)
could enhance the radiated signal, effectively replacing a planetary-mass
modulated source with a laboratory-scale resonant emitter.

\subsection{Why $\Qpsi$ Does NOT Help: The Radiation Problem}

\begin{theorem}[$\Qpsi$ Does Not Reduce Required Transmitter Mass]
\label{thm:Q-no-help}
The far-field psi-radiation from a resonant cavity is independent of $\Qpsi$
at fixed input power.

\begin{proof}
The radiated power from a psi-cavity of stored energy $E$ and quality factor
$\Qpsi$ is:
\begin{equation}
P_{\mathrm{rad}} = \frac{\Opsi \cdot E_{\mathrm{stored}}}{\Qpsi}
= \frac{\Opsi}{\Qpsi} \cdot \frac{\Qpsi}{\Opsi} \cdot P_{\mathrm{in}}
= P_{\mathrm{in}}.
\label{eq:P-rad}
\end{equation}
In steady state, the radiated power equals the input power. The $\Qpsi$
factor cancels exactly. This is a universal property of driven resonant
systems: high $Q$ increases stored energy but \emph{not} radiated power
in steady state.

The psi-field amplitude at distance $d$ from a radiating cavity:
\begin{equation}
\delta\psi(d) \sim \frac{\GN P_{\mathrm{in}}}{c^5 d}
= \frac{6.67\ee{-11} \cdot P_{\mathrm{in}}}{2.43\ee{42} \cdot d}
= \frac{2.74\ee{-53} \cdot P_{\mathrm{in}}}{d}.
\label{eq:psi-rad-field}
\end{equation}
This depends only on $P_{\mathrm{in}}$ and $d$, not on $\Qpsi$.
\end{proof}
\end{theorem}

\begin{finding}[$\Qpsi$ Amplification: Irrelevant for Psi-Communications]
\label{find:Q-irrelevant}
\killed{Resonant cavity $\Qpsi$ amplification does NOT reduce the
transmitter mass or power requirements for psi-communications.}

\textbf{Physical intuition}: A high-$Q$ cavity is like a very resonant
bell --- it rings for a long time after being struck, but it doesn't
radiate more total sound energy than was put in. The $\Qpsi = 10^9$
factor means the cavity \emph{stores} a lot of psi-energy, but it
\emph{radiates} at the same rate as the input power.

\textbf{Distinction from EM cavities}: In EM communications, a high-$Q$
antenna can be efficient (radiating most of the input power as EM waves).
But psi-radiation from laboratory EM sources is suppressed by $\GN/c^5$,
and no cavity $Q$ can overcome this fundamental coupling constant.
\end{finding}

\subsection{Other Modulation Mechanisms Without Planetary Mass}

We systematically evaluate every conceivable mechanism:

\begin{table}[H]
\centering
\caption{Alternative psi-modulation mechanisms: can any avoid planetary mass?}
\label{tab:P8-mechanisms}
\begin{tabular}{@{}p{3cm}p{4.5cm}p{5cm}@{}}
\toprule
Mechanism & Physics & Verdict \\
\midrule
Mass modulation (baseline P8)
  & $\delta\psi = GM_s\xi_0/(c^2 d)$
  & Requires $M_s \sim 10^{10}$~kg at 500~m (W3i6) \\[3pt]
Resonant cavity ($\Qpsi$)
  & $P_{\mathrm{rad}} = P_{\mathrm{in}}$ (independent of $Q$)
  & \killed{No improvement} (Theorem~\ref{thm:Q-no-help}) \\[3pt]
EM-to-$\psi$ (Process A)
  & $\delta\psi \propto (\lambda-1) \cdot u_{\mathrm{EM}} \cdot \GN/c^4$
  & Requires $(\lambda-1)$; even at $\lambdaone = 1$, coupling
    $\GN/c^4$ kills it. P15 transmitter already killed (W3i5). \\[3pt]
Coherent enhancement (superradiance)
  & $N_{\mathrm{cav}}$ synchronised cavities: $P_{\mathrm{rad}}
    \propto N_{\mathrm{cav}}^2 P_1$ (if coherent)
  & Superradiant scaling requires $N_{\mathrm{cav}}$ cavities within
    $\lambda_\psi/2\pi$. At $f = 1$~GHz: $\lambda_\psi = 0.3$~m.
    Packing $N$ cavities in 0.3~m is impractical for $N > 1$. \\[3pt]
Nuclear energy density
  & Modulate nuclear reactions: $u_{\mathrm{nuc}} \sim 10^{17}$~J/m$^3$
  & Nuclear energy is rest mass ($E = mc^2$), already captured by
    $\psigrav$. No enhancement over mass modulation. \\[3pt]
Rotating mass (Weber bar)
  & Quadrupole radiation: $P \sim G M^2 R^4 \omega^6/c^5$
  & For $M = 10^6$~kg, $R = 1$~m, $\omega = 100$~rad/s:
    $P \sim 10^{-26}$~W. Negligible. \\
\bottomrule
\end{tabular}
\end{table}

\begin{finding}[No Modulation Mechanism Avoids Planetary Mass]
\label{find:no-mechanism}
\killed{Every conceivable psi-modulation mechanism encounters the same
fundamental barrier: the $\GN/c^4$ or $\GN/c^5$ coupling factor that
suppresses matter-to-$\psi$ conversion. No resonant, coherent,
electromagnetic, or nuclear mechanism can overcome this by more than
a few orders of magnitude (at most $10^2$--$10^3$), leaving a
$>10^{8}\times$ gap to useful transmitter performance.}
\end{finding}

%=============================================================================
\section{Minimum Transmitter Mass for 1~kbit/s at 1~AU}
\label{sec:P8-minimum-mass}
%=============================================================================

\subsection{Link Budget Derivation}

The psi-channel capacity from a mass-modulated source (W3i4/W3i6):
\begin{equation}
C = B\log_2\!\left(1 + \SNR\right),
\label{eq:capacity}
\end{equation}
where:
\begin{equation}
\SNR = \frac{S_{\mathrm{signal}}}{S_{\mathrm{noise}}}
= \frac{(2\GN M_s \xi_0)^2}{c^4 d^2 \cdot S_h(f) \cdot B},
\label{eq:SNR-link}
\end{equation}
with:
\begin{itemize}[nosep]
\item $M_s$ = source mass (kg), modulation depth $\xi_0$ (dimensionless,
  $\sim 10^{-3}$ for mechanical modulation),
\item $d$ = distance = 1~AU $= 1.5\ee{11}$~m,
\item $B$ = bandwidth (Hz); we take $B = C$ (Shannon limit),
\item $S_h(f)$ = receiver noise PSD; at the RPN limit for $N$ sensors:
  $S_h^{1/2} = 2.94\ee{-15}/\sqrt{N}\sqrtHz$.
\end{itemize}

For $C = 1$~kbit/s, $B = 1$~kHz:
\begin{equation}
\SNR_{\mathrm{req}} = 2^{C/B} - 1 = 2^1 - 1 = 1.
\label{eq:SNR-req}
\end{equation}

Setting $\SNR = 1$ and solving for $M_s$:
\begin{equation}
M_s = \frac{c^2 d}{2\GN \xi_0}
\sqrt{S_h(f) \cdot B}.
\label{eq:M-min}
\end{equation}

\subsection{Numerical Evaluation}

With $N = 10^6$ sensors (best conceivable array):
\begin{equation}
S_h^{1/2} = \frac{2.94\ee{-15}}{10^3} = 2.94\ee{-18}\sqrtHz.
\end{equation}

\begin{align}
M_s &= \frac{(3\ee{8})^2 \times 1.5\ee{11}}{2 \times 6.67\ee{-11} \times 10^{-3}}
\times \sqrt{(2.94\ee{-18})^2 \times 10^3} \nonumber\\
&= \frac{9\ee{16} \times 1.5\ee{11}}{1.33\ee{-13}}
\times \sqrt{8.64\ee{-33}} \nonumber\\
&= \frac{1.35\ee{28}}{1.33\ee{-13}}
\times 9.3\ee{-17} \nonumber\\
&= 1.015\ee{41} \times 9.3\ee{-17} \nonumber\\
&= 9.4\ee{24}\,\mathrm{kg}.
\label{eq:M-min-numerical}
\end{align}

This is $\sim 1.6\,M_\oplus$ (1.6 Earth masses).

\begin{correction}[Minimum Mass at 1~AU Depends Critically on Receiver Array]
\label{corr:mass-1AU}
The minimum transmitter mass for 1~kbit/s at 1~AU depends on $N$:

\begin{center}
\begin{tabular}{@{}lcc@{}}
\toprule
Receiver array $N$ & $S_h^{1/2}$ ($\sqrtHz$) & $M_s^{\min}$ (kg) \\
\midrule
$10^0$ (single sensor) & $2.94\ee{-15}$ & $3.0\ee{27}$ \\
$10^2$ & $2.94\ee{-16}$ & $3.0\ee{26}$ \\
$10^4$ & $2.94\ee{-17}$ & $3.0\ee{25}$ \\
$10^6$ & $2.94\ee{-18}$ & $9.4\ee{24}$ \\
\bottomrule
\end{tabular}
\end{center}

\textbf{All values are planetary-mass or above.} Even with $10^6$ sensors,
the minimum mass is 1.6~Earth masses.
\end{correction}

\subsection{At Shorter Range: Submarine Application Revisited}

For the submarine killer application ($d = 500$~m, $C = 1$~kbit/s,
$N = 10^4$ receiver array):
\begin{align}
M_s &= \frac{(3\ee{8})^2 \times 500}{2 \times 6.67\ee{-11} \times 10^{-3}}
\times \sqrt{(2.94\ee{-17})^2 \times 10^3} \nonumber\\
&= \frac{4.5\ee{19}}{1.33\ee{-13}} \times 9.3\ee{-16} \nonumber\\
&= 3.38\ee{32} \times 9.3\ee{-16} \nonumber\\
&= 3.1\ee{17}\,\mathrm{kg}.
\label{eq:M-sub-corrected}
\end{align}

\begin{finding}[Corrected Submarine Transmitter Mass]
\label{find:sub-mass}
For 1~kbit/s at 500~m with $N = 10^4$ sensors, the minimum transmitter
mass is $M_s \approx 3\ee{17}$~kg $\approx 5\ee{-8}\,\Msun$. This is
roughly the mass of a large asteroid or small dwarf planet ($\sim$Ceres).

The W3i6 estimate of $M_s \sim 10^{10}$~kg assumed a \emph{less sensitive}
receiver (higher noise floor). With the corrected RPN-limited P6 receiver
at $N = 10^4$, the mass requirement is \emph{higher} because the W3i6
calculation used a different noise model.

\honest{The discrepancy with W3i6 ($10^{10}$~kg vs $3\ee{17}$~kg)
arises from the different noise models used. The W3i6 calculation
(Eq.~8 in W3i6\_P8\_P11\_update.tex) used a thermal noise floor
$N_0 = 2.5\ee{-38}$, which is far below the quantum noise limit. With
the correct quantum-limited P6 receiver noise, the mass requirement is
much higher. This \textbf{worsens} the feasibility assessment.}
\end{finding}

\begin{tcolorbox}[colback=red!5!white, colframe=red!60!black,
  title=\textbf{P8 W3i7 Verdict: Mass Gap Worse Than Previously Assessed}]
\begin{itemize}[nosep]
\item Resonant cavity $\Qpsi$: does NOT help (Theorem~\ref{thm:Q-no-help}).
\item No alternative modulation mechanism avoids planetary mass.
\item Minimum mass for 1~kbit/s at 1~AU: $\sim 10^{25}$~kg (1.6~$M_\oplus$).
\item Minimum mass for 1~kbit/s at 500~m (submarine): $\sim 3\ee{17}$~kg
  (Ceres-mass).
\item Mass gap from achievable technology ($\sim 10^6$~kg): $\sim 10^{11}\times$
  for submarine, $\sim 10^{19}\times$ for interplanetary.
\item \killed{Psi-communications is even more infeasible than W3i6 assessed,
  when using quantum-limited receiver noise.}
\end{itemize}
\end{tcolorbox}

%=============================================================================
\part{Patent 9: SQL vs Heisenberg Positioning --- Application Survival}
\label{part:P9}
%=============================================================================

%=============================================================================
\section{The Heisenberg Limit vs Standard Quantum Limit}
\label{sec:P9-HL-SQL}
%=============================================================================

\subsection{Definitions}

For a $\psi$-field gradient measurement using $N_a$ atoms in an atom
interferometer with interrogation time $T$:

\begin{definition}[Standard Quantum Limit (SQL)]
\label{def:SQL}
The SQL arises from uncorrelated quantum projection noise:
\begin{equation}
\sigma_\psi^{\mathrm{SQL}} = \frac{1}{k_{\mathrm{eff}} T^2 \sqrt{N_a}},
\label{eq:SQL}
\end{equation}
where $k_{\mathrm{eff}}$ is the effective wave vector and $N_a$ is the
atom number per shot.
\end{definition}

\begin{definition}[Heisenberg Limit (HL)]
\label{def:HL}
The HL uses maximally entangled states (GHZ, spin-squeezed, or NOON states)
to achieve:
\begin{equation}
\sigma_\psi^{\mathrm{HL}} = \frac{1}{k_{\mathrm{eff}} T^2 N_a},
\label{eq:HL}
\end{equation}
which is $\sqrt{N_a}$ times better than the SQL.
\end{definition}

The ratio:
\begin{equation}
\frac{\sigma^{\mathrm{SQL}}}{\sigma^{\mathrm{HL}}} = \sqrt{N_a}.
\label{eq:SQL-HL-ratio}
\end{equation}

\subsection{Typical $N_a$ Values}

For state-of-the-art atom interferometers:
\begin{itemize}[nosep]
\item Current laboratory: $N_a \sim 10^6$ atoms per shot.
\item Projected (MAGIS/AION class): $N_a \sim 10^8$.
\item Far-future: $N_a \sim 10^{10}$.
\end{itemize}

The SQL-to-HL improvement factor:
\begin{itemize}[nosep]
\item $N_a = 10^6$: $\sqrt{N_a} = 10^3$ (i.e., SQL is $10^3\times$ worse).
\item $N_a = 10^8$: $\sqrt{N_a} = 10^4$.
\item $N_a = 10^{10}$: $\sqrt{N_a} = 10^5$.
\end{itemize}

%=============================================================================
\section{P9 Positioning: SQL vs HL Performance}
\label{sec:P9-SQL-performance}
%=============================================================================

\subsection{Single-Sensor Positioning}

The W3i4/W3i6 single-sensor positioning precision at the SQL:
\begin{equation}
\sigma_r^{\mathrm{SQL}} = \frac{c^2}{g} \cdot \sigma_\psi^{\mathrm{SQL}}
= \frac{c^2}{g} \cdot \frac{1}{k_{\mathrm{eff}} T^2 \sqrt{N_a}}.
\label{eq:sigma-r-SQL-single}
\end{equation}

From W3i3/W3i4, the single-sensor precision is $\sigma_r = 0.091$~mm
at the SQL (this was always SQL, since entanglement was not assumed for
single sensors).

\confirmed{The 0.091~mm single-sensor precision is already at the SQL.
No correction needed for single-sensor operation.}

\subsection{Array Positioning: SQL vs HL}

The W3i4/W3i6 array positioning precision is:
\begin{equation}
\sigma_r^{\mathrm{array}} = \frac{\sigma_r^{\mathrm{single}}}{\sqrt{N_s} \cdot \eta},
\label{eq:sigma-r-array}
\end{equation}
where $N_s$ is the number of sensors in the array and $\eta$ is the
entanglement enhancement factor:
\begin{itemize}[nosep]
\item $\eta = 1$ at SQL (no entanglement).
\item $\eta = \sqrt{N_a}$ at HL (full Heisenberg scaling \emph{within each sensor}).
\end{itemize}

The 8.7~nm claim from W3i4:
\begin{equation}
\sigma_r^{\mathrm{HL}} = \frac{0.091\,\mathrm{mm}}{\sqrt{N_s} \cdot \sqrt{N_a}}.
\end{equation}
For $N_s = 100$ sensors and $N_a = 10^6$ atoms:
\begin{equation}
\sigma_r^{\mathrm{HL}} = \frac{9.1\ee{-5}}{10 \times 10^3}
= \frac{9.1\ee{-5}}{10^4} = 9.1\ee{-9}\,\mathrm{m} = 9.1\,\mathrm{nm}.
\label{eq:HL-array}
\end{equation}
This confirms the 8.7~nm claim (within rounding).

At the SQL (no intra-sensor entanglement, $\eta = 1$):
\begin{equation}
\sigma_r^{\mathrm{SQL}} = \frac{0.091\,\mathrm{mm}}{\sqrt{100}}
= \frac{9.1\ee{-5}}{10} = 9.1\ee{-6}\,\mathrm{m}
= 9.1\,\mu\mathrm{m}.
\label{eq:SQL-array}
\end{equation}

\begin{finding}[SQL Array Positioning: $\sim 9~\mu$m]
\label{find:SQL-positioning}
\corrected{At the SQL (no entanglement), a 100-sensor P9 array achieves
$\sigma_r^{\mathrm{SQL}} \approx 9~\mu$m, NOT 9~nm.}

The $10^3\times$ degradation from HL to SQL reflects the $\sqrt{N_a}$
entanglement advantage at $N_a = 10^6$.

\textbf{State of the art for Heisenberg-limited sensors}: As of 2025,
the best demonstrated spin-squeezing in atom interferometers achieves
$\sim 10$--$20$~dB below the SQL ($\eta \sim 3$--$10$), far from the
Heisenberg limit ($\eta = \sqrt{N_a} \sim 10^3$). The HL claim requires
a $\sim 100\times$ improvement over current squeezing technology.
\end{finding}

\subsection{Intermediate Squeezing Scenarios}

\begin{table}[H]
\centering
\caption{P9 array positioning precision ($N_s = 100$ sensors,
$N_a = 10^6$ atoms) as a function of squeezing level.}
\label{tab:P9-squeezing}
\begin{tabular}{@{}lccc@{}}
\toprule
Squeezing Level & $\eta$ & $\sigma_r$ (array) & Technology Status \\
\midrule
SQL (no squeezing) & 1 & $9.1\,\mu$m & Available now \\
10~dB squeezing & 3.2 & $2.8\,\mu$m & State of art (2025) \\
20~dB squeezing & 10 & $0.91\,\mu$m & Near-term ($\sim$2030) \\
30~dB squeezing & 31.6 & $0.29\,\mu$m & Challenging \\
40~dB squeezing & 100 & $91$~nm & Very challenging \\
50~dB squeezing & 316 & $29$~nm & Frontier \\
Heisenberg limit (60~dB) & 1000 & $9.1$~nm & Aspirational \\
\bottomrule
\end{tabular}
\end{table}

%=============================================================================
\section{Application Survival Analysis: SQL vs HL}
\label{sec:P9-apps}
%=============================================================================

\subsection{Application Requirements}

\begin{table}[H]
\centering
\caption{P9 applications and their positioning precision requirements.
``Survives at SQL'' means the SQL precision ($9~\mu$m array, $0.09$~mm
single sensor) meets the application requirement.}
\label{tab:P9-apps}
\begin{tabular}{@{}p{4cm}cccl@{}}
\toprule
Application & Required $\sigma_r$ & SQL OK? & HL needed?
  & Notes \\
\midrule
\multicolumn{5}{@{}l}{\textbf{SURVIVE AT SQL ($\sigma_r \leq 9~\mu$m array, 0.09~mm single):}} \\[3pt]
Submarine navigation & $< 1$~m & \textbf{Yes} & No
  & Single sensor sufficient \\
Underwater positioning & $< 10$~cm & \textbf{Yes} & No
  & Single sensor sufficient \\
Geological survey (ore body) & $< 1$~m & \textbf{Yes} & No
  & Resolution is depth-limited \\
Seismic monitoring & $< 1$~mm & \textbf{Yes} & No
  & Single sensor sufficient \\
Autonomous vehicle nav & $< 10$~cm & \textbf{Yes} & No
  & Single sensor sufficient \\
Tunnel/mine surveying & $< 1$~cm & \textbf{Yes} & No
  & Single sensor sufficient \\
Infrastructure monitoring & $< 1$~mm & \textbf{Yes} & No
  & Single sensor sufficient \\
Geoid mapping & $< 1$~cm & \textbf{Yes} & No
  & Array improves coverage \\
Tectonic plate monitoring & $< 1$~mm & \textbf{Yes} & No
  & Array for long baselines \\[6pt]
\midrule
\multicolumn{5}{@{}l}{\textbf{REQUIRE SQUEEZING ($\sigma_r < 9~\mu$m):}} \\[3pt]
Precision geodesy (VLBI-class) & $\sim 1~\mu$m & 20~dB & No
  & Competitive with VLBI \\
Gravitational gradiometry & $\sim 0.1~\mu$m & 40~dB & No
  & Sub-$\mu$m gradient maps \\[6pt]
\midrule
\multicolumn{5}{@{}l}{\textbf{REQUIRE HEISENBERG LIMIT ($\sigma_r < 100$~nm):}} \\[3pt]
GW detector calibration & $\sim 10$~nm & \textbf{No} & \textbf{Yes}
  & Independent verification of LISA \\
Interferometric geodesy & $\sim 1$~nm & \textbf{No} & \textbf{Yes}
  & Fundamental reference frame \\
Chip-scale metrology & $< 100$~nm & \textbf{No} & \textbf{Yes}
  & Semiconductor positioning \\
\bottomrule
\end{tabular}
\end{table}

\begin{finding}[Most P9 Applications Survive at SQL]
\label{find:P9-SQL-survive}
\verdict{The dominant practical applications of P9 --- submarine navigation,
underwater positioning, geological survey, seismic monitoring, autonomous
vehicle navigation, tunnel surveying, and infrastructure monitoring ---
ALL survive at the SQL with wide margin.}

The single-sensor precision of $0.091$~mm already exceeds the requirements
of every major application by factors of $10\times$--$10{,}000\times$.
The 8.7~nm Heisenberg-limited array precision is relevant only for:
\begin{enumerate}[nosep]
\item GW detector calibration (independent nm-level distance verification).
\item Fundamental geodesy (Earth shape at nm precision).
\item Semiconductor positioning (requires competitive nm accuracy).
\end{enumerate}
These are important but \emph{niche} applications compared to the
broad utility of sub-mm positioning.
\end{finding}

%=============================================================================
\section{Underwater Positioning: Detailed SQL Analysis}
\label{sec:P9-underwater}
%=============================================================================

\subsection{The Underwater Environment}

The P9 underwater positioning system operates by measuring $\nabla\psigrav$
from the seafloor, coastline, and geological features. The key advantage
over acoustic or EM positioning is that the $\psi$-field is unaffected by
water (confirmed W3i6: $\dpsiem$ noise $< 10^{-25}$ relative to signal).

\begin{finding}[Underwater Positioning at SQL: World-Class Performance]
\label{find:P9-underwater-SQL}
At the SQL, a single P9 sensor provides $\sigma_r = 0.091$~mm underwater
positioning. For comparison:

\begin{center}
\begin{tabular}{@{}lcc@{}}
\toprule
Technology & Precision & Range \\
\midrule
GPS (surface) & $\sim 1$~m & Surface only \\
USBL acoustic & $\sim 0.1$--$1$~m & $< 10$~km \\
LBL acoustic & $\sim 1$--$10$~cm & $< 5$~km (array) \\
INS (unaided, 1~hr) & $\sim 1$~km drift & Unlimited \\
INS + Doppler velocity & $\sim 10$~m (1~hr) & Unlimited \\
\textbf{DFD P9 (SQL, single)} & $\mathbf{0.091}$~mm & Unlimited \\
\textbf{DFD P9 (SQL, 100-array)} & $\mathbf{9~\mu}$m & Unlimited \\
\textbf{DFD P9 (HL, 100-array)} & $\mathbf{9}$~nm & Unlimited \\
\bottomrule
\end{tabular}
\end{center}

Even at the SQL, the DFD P9 single sensor achieves $\sim 10{,}000\times$
better precision than the best acoustic positioning (LBL at 1~cm). The
Heisenberg limit is $\sim 10^7\times$ better than LBL, but the SQL
advantage is already transformative.
\end{finding}

%=============================================================================
\section{Geological Detection: SQL vs HL Impact}
\label{sec:P9-geo-SQL}
%=============================================================================

\subsection{Minimum Detectable Anomaly Mass}

From W3i6 (confirmed), the minimum detectable anomaly mass at depth $D$
with sensor noise $\sigma_\psi$:
\begin{equation}
M_{\mathrm{anom}}^{\min} = \frac{c^2 D^2 \sigma_\psi}{G \cdot \mathrm{SNR}_{\min}}.
\label{eq:M-anom-min}
\end{equation}

At 23~km depth with SNR $= 5$:

\begin{table}[H]
\centering
\caption{Minimum detectable anomaly mass at 23~km depth, SQL vs HL.}
\label{tab:P9-geo-SQL}
\begin{tabular}{@{}lccc@{}}
\toprule
Sensor Configuration & $\sigma_\psi$ & $M_{\mathrm{anom}}^{\min}$ (kg)
  & Equivalent \\
\midrule
Single sensor (SQL) & $\sim 10^{-15}$ & $1.56\ee{11}$ & $\sim 0.1$~km$^3$ rock \\
100-sensor array (SQL) & $\sim 10^{-16}$ & $1.56\ee{10}$ & $\sim 10^4$~m$^3$ rock \\
100-sensor array (HL) & $\sim 10^{-19}$ & $1.56\ee{7}$ & $\sim 10$~m$^3$ rock \\
$10^6$-sensor array (HL) & $\sim 10^{-22}$ & $1.56\ee{4}$ & $\sim 10$~L rock \\
\bottomrule
\end{tabular}
\end{table}

\begin{finding}[Geological Detection: SQL Sufficient for Ore Body Discovery]
\label{find:P9-geo-SQL}
At the SQL with a 100-sensor array, P9 can detect geological anomalies of
$\sim 1.6\ee{10}$~kg ($\sim 6000$~m$^3$ of rock density contrast) at 23~km
depth. This corresponds to a $\sim 18$~m diameter ore body --- well within
the range of commercially valuable mineral deposits (typical copper/gold
ore bodies are $>100$~m diameter).

\verdict{The SQL is sufficient for geological survey and mineral
exploration. The Heisenberg limit enables detection of $\sim 10$~m$^3$
anomalies, which would be relevant for archaeological surveys, tunnel
detection, and cavity mapping, but is not needed for mineral exploration.}
\end{finding}

%=============================================================================
\part{Cross-Patent Summary and Open Items}
\label{part:summary}
%=============================================================================

%=============================================================================
\section{Definitive Status After W3i7}
\label{sec:status}
%=============================================================================

\begin{tcolorbox}[colback=green!5!white, colframe=green!40!black,
  title=\textbf{P6: Decihertz Detection --- Context-Dependent Claim}]
\begin{itemize}[nosep]
\item DFD P6 does NOT beat LISA, DECIGO, BBO, ET, or CE on raw sensitivity
  at any frequency --- \textbf{New W3i7 finding}.
\item DFD P6 is uniquely competitive as a \emph{ground-based decihertz
  detector} in the 0.01--1~Hz gap where no funded mission operates ---
  \textbf{New W3i7 finding}.
\item Only detectable source at feasible $N$: IMBH binary merger at
  $N \sim 10^6$ --- \textbf{Confirmed from W3i4}.
\item Technology synergy with MAGIS/AION makes incremental deployment
  realistic --- \textbf{New W3i7 assessment}.
\item ``Decihertz GW detection'' claim is valid but must be qualified:
  it fills a gap, it does not outperform --- \textbf{New W3i7 verdict}.
\end{itemize}
\end{tcolorbox}

\begin{tcolorbox}[colback=red!5!white, colframe=red!60!black,
  title=\textbf{P8: Communications --- Even More Infeasible}]
\begin{itemize}[nosep]
\item $\Qpsi$ amplification: does NOT help (radiated power = input power,
  independent of $Q$) --- \textbf{New W3i7 theorem}.
\item No alternative modulation mechanism avoids planetary mass ---
  \textbf{New W3i7 finding}.
\item Minimum mass for 1~kbit/s at 1~AU ($N = 10^6$ receiver):
  $\sim 10^{25}$~kg ($\sim 1.6\,M_\oplus$) --- \textbf{New W3i7 result}.
\item Minimum mass for 1~kbit/s at 500~m ($N = 10^4$ receiver):
  $\sim 3\ee{17}$~kg (Ceres-mass) --- \textbf{New W3i7 result}.
\item Mass gap from technology: $10^{11}\times$ (submarine),
  $10^{19}\times$ (interplanetary) --- \textbf{New W3i7 assessment}.
\item \killed{Psi-communications feasibility is WORSE than W3i6 assessed
  when using quantum-limited noise.}
\end{itemize}
\end{tcolorbox}

\begin{tcolorbox}[colback=green!5!white, colframe=green!40!black,
  title=\textbf{P9: Positioning --- SQL Sufficient for Major Applications}]
\begin{itemize}[nosep]
\item 8.7~nm array positioning requires Heisenberg-limited sensors ---
  \textbf{Confirmed}.
\item SQL array positioning: $\sim 9~\mu$m (100-sensor array,
  $N_a = 10^6$) --- \textbf{New W3i7 result}.
\item Single-sensor SQL: 0.091~mm --- \textbf{Confirmed from W3i4}.
\item All major applications (submarine nav, underwater positioning,
  geological survey, seismic monitoring, autonomous vehicles) survive
  at SQL --- \textbf{New W3i7 finding}.
\item HL needed only for: GW calibration, fundamental geodesy,
  semiconductor metrology --- \textbf{New W3i7 finding}.
\item Underwater positioning at SQL: $\sim 10{,}000\times$ better than
  best acoustic (LBL) --- \textbf{New W3i7 assessment}.
\item Geological survey at SQL: detects $\sim 18$~m diameter ore bodies
  at 23~km depth --- \textbf{New W3i7 result}.
\end{itemize}
\end{tcolorbox}

%=============================================================================
\section{Open Items for W3i8}
\label{sec:open}
%=============================================================================

\begin{enumerate}
\item \textbf{P6}: The MAGIS-100 pathfinder (expected $\sim$2027) will
  provide the first experimental constraint on DFD P6-type atom
  interferometer sensitivity in the decihertz band. What is the expected
  MAGIS-100 sensitivity in DFD P6 terms?

\item \textbf{P8}: With psi-comms definitively infeasible by
  $> 10^{11}\times$, should P8 be formally retired from the active patent
  portfolio, or does the theoretical uniqueness (only physics penetrating
  arbitrary matter) preserve long-term value?

\item \textbf{P9}: Current squeezing technology achieves $\sim 10$--$20$~dB
  (corresponding to $\sigma_r \sim 1$--$3~\mu$m). What is the realistic
  roadmap for 30--40~dB squeezing, and what applications become accessible
  at each level?

\item \textbf{Cross-patent}: The P6 receiver noise floor determines both
  GW detection sensitivity (P6) and communication channel capacity (P8).
  Any improvement in P6 sensor technology simultaneously improves both,
  but the P8 mass gap of $10^{11}\times$ means sensor improvements alone
  cannot rescue P8.
\end{enumerate}

\end{document}
